\begin{resumo}
\begin{SingleSpace}
Este trabalho investiga em que medida o tipo e a ordem das perspectivas atribuídas a agentes baseados em grandes modelos de linguagem alteram o conteúdo e o tom afetivo de notícias ao longo de cadeias de reescrita. Para mensurar essas transformações, propõe o \textit{Structured Topic Drift Index} (STDI), que decompõe a deriva em tema, subtópicos, entidades e relações, com incrementos de contradição interna e mudança afetiva. O estudo reúne seis rodadas finais, combinando três modelos geradores e dois corpora de 50 entradas cada: descrições de notícias políticas em inglês e legendas de publicações jornalísticas em português no contexto eleitoral brasileiro. Cada rodada aplica oito trajetórias de quatro posições, organizadas em árvores com etapas iniciais compartilhadas, totalizando 7.200 reescritas únicas. A análise combina comparações quantitativas pareadas por entrada, inspeções qualitativas e estimativas de valência, ativação e dominância, distinguindo afastamento do original e mudança incremental. Os resultados mostram que as reescritas podem preservar o acontecimento central e seus participantes, mas alterar justificativas, motivos e relações. Subtópicos e relações apresentam as maiores médias dos componentes de conteúdo. A retransmissão neutra desde o início está associada a maior preservação, enquanto passagens neutras posteriores conservam interpretações recebidas sem recuperar automaticamente o original. Trajetórias com perspectiva conspiratória apresentam, em média, menor valência e maior ativação segundo o avaliador afetivo por modelo de linguagem. As comparações não identificam uma ordenação que produza maior afastamento de forma consistente entre geradores e corpora. Divergências entre métodos e oscilações de escores em textos equivalentes também revelam limites da mensuração. O estudo conclui que a análise de trajetórias e componentes permite localizar transformações que um escore final isolado não explica. As contribuições são instrumentais e descritivas, limitadas pela seleção dos corpora, pelo avaliador comum e pela ausência de referência humana numérica independente. O STDI descreve mudanças textuais, sem estabelecer veracidade factual externa ou comportamento humano de circulação de notícias.
\end{SingleSpace}
\vspace{\onelineskip}
\textbf{Palavras-chave:} deriva informacional; inteligência artificial (AI); grandes modelos de linguagem (LLM); agentes computacionais; reescrita de notícias; análise afetiva.

\end{resumo}
